\section{How Part 10 was run}\label{app:provenance}

Part 10 and this paper were produced by an AI system, Claude (Anthropic), on 27--28 October 2026. Kaarel H\"anni posed the request: ``a couple more hours'' on a better discrete learner, then interpretability of the best learner so far, and then a paper about it that explains what kind of learner it is, how well it does against the baselines, and what the interpretability shows. No human has reviewed the code, the results or this text. Times are UTC. The records are:
\begin{itemize}
\item the protocol \repo{study/part10/PROTOCOL.md}, written before any design work and amended only by dated decision entries and notes;
\item the builders' notes \repo{study/part10/p10_sum/DESIGN.md}, \repo{study/part10/p10_hier/DESIGN.md} and \repo{study/part10/p10_codes/DESIGN.md};
\item the assembly record \repo{study/part10/hyb_portfolio3/ASSEMBLY.md};
\item the verification record of the interpretability analyses, \repo{study/part10/interp/VERIFICATION.json}.
\end{itemize}

\paragraph{Organisation.} One orchestrating session wrote the protocol, fixed the held-out sets, ran every held-out evaluation and reference run, applied the rules and assembled v3. It delegated five kinds of work to separate agents with their own instructions: the three builders; two interpretability analysts; four verifiers of the analyses; the drafting of sections of this paper; and three checkers of its claims (\cref{app:repro}). The builders worked on development data only. The orchestrator read no held-out result of a Part 10 learner before that learner was frozen (\cref{sec:disc-provenance-disclosures} lists what it did read).

\subsection{Protocol and builders}\label{app:builders}

\paragraph{The protocol} was committed at 18:50 on 27 October (\texttt{dc70266c}), before any design work. It fixed the held-out sets \code{dlheld7}, \code{qd_c} and \code{qd_d} and the development draw \code{dlfamdev}; the budgets; the three builders' tasks and their fixed development evaluations; the integration rule; the held-out evaluation with its references; the best-learner rule; and the plan for the interpretability analysis and this paper (\cref{sec:discipline}).

\begin{table}[ht]
\centering\footnotesize\setlength{\tabcolsep}{4pt}\renewcommand{\arraystretch}{1.15}
\begin{tabularx}{\linewidth}{>{\raggedright\arraybackslash}p{2.0cm}>{\raggedright\arraybackslash}X>{\raggedright\arraybackslash}p{2.1cm}>{\raggedright\arraybackslash}p{2.0cm}>{\raggedright\arraybackslash}p{3.6cm}}
\toprule
builder & component & development & frozen (md5) & fixed development evaluation\\
\midrule
T, \code{p10_sum} & DFG v1.3: weighted sum blocks, forced blocks at $n\ge48$, beam at $n\le128$ & 18:52--19:44 (52 min) & 19:44 (\texttt{5ee995ad}; \texttt{dfg13.py} \texttt{69cabe5c}) & 20:00--21:35 on \code{designed}, \code{dlsymdev}, \code{dlfamdev}, 3 seeds; kr\_vs\_kp at $n=1000$\\
I1, \code{p10_hier} & a max-pooling second level in \code{ggbt}'s bank; pairs dropped & 18:52--19:45 (53 min) & 19:45 (\texttt{1da21bf7}) & 19:45--21:35 on the 8 image cells and \code{qd_a} at $n=60000$, seed 0\\
I2, \code{p10_codes} & a second booster \code{cgbt} on pooled sums and bars & 21:37--22:10 (33 min) & 22:10 (\texttt{fd688acd}) & 22:10--23:43, the same 9 cells; interrupted at 22:40 and resumed (below)\\
\bottomrule
\end{tabularx}
\caption{The three Part 10 builders, 27 October. Times from each builder's \texttt{DESIGN.md}; the evaluation's end for \code{p10_codes} is the commit that recorded it (\texttt{16ef5fde}). Every evaluation cell records the md5 of the file it ran.}
\label{tab:builders}
\end{table}

\paragraph{The builders} are summarized in \cref{tab:builders}. Each changed one component of v2, delivered a full learner file, a \texttt{DESIGN.md} and an md5 at its freeze, and then ran the fixed development evaluation once. The image builders used validation rows of the development sets only. The toy builder scored its DFG-alone checks on fresh validation draws, and its portfolio checks on the toy harness's test points of the development suites (\code{designed}, \code{dlsymdev}, \code{dlfamdev}, \code{dlheld6}). No builder used a held-out set. Three events during the evaluations are recorded in the notes:
\begin{itemize}
\item \emph{\code{p10_sum}.} Its evaluation was started at 19:44 as 8 shards in one shell session. The machine's autogroup scheduler gave the shards together about 0.55 cores, so at 20:00 they were stopped, before any result was written, and restarted unchanged in four sessions. Only the restarted run produced results.
\item \emph{\code{p10_hier}.} Its \code{qd_dev} process was killed by the shared memory limit at about 20:29 and restarted; the finished $n=1000$ cell was kept and $n=10000$ refitted with the same random streams. Its $n=60000$ cell waited 1 h 28 min for the shared large-fit lock and ran 21:13--21:35.
\item \emph{\code{p10_codes}.} It started after the other two builders had finished. The container restarted at 22:40, with 2 of its 9 cells done, and the workflow that ran the builders was lost. The orchestrator resumed the evaluation with the builder's own script. The cells are resumable and record the learner's md5, which was unchanged. The orchestrator took the builder's summary from its \texttt{DESIGN.md}, whose evaluation and cost placeholders (\S3, \S5) the builder never filled.
\end{itemize}

\subsection{The integration and best-learner decisions}\label{sec:decisions}

The rules (\cref{sec:discipline}) were applied by \repo{study/part10/compare_part10.py} from the result files, and each decision was entered in the protocol when it was made. The numbers are in \repo{study/part10/RESULTS.md}.
\begin{itemize}
\item \emph{Toy component: adopted} (21:35; commit \texttt{33df158e}, 21:50). Against v2 on \code{designed}, \code{dlsymdev} and \code{dlfamdev} (41 targets, inputs permuted, 3 seeds), \code{p10_sum}'s mean $\log_2\nstar$ ratio is $-0.066$ over the 40 targets both learn: $-0.028$ on \code{designed}, $-0.178$ on \code{dlsymdev} and $-0.055$ on \code{dlfamdev}. The largest increase is $+0.26$ (fd\_06). The rule asked for a mean below $-0.05$ and no increase above 0.5. v2's development fits are \code{p9_sum}'s on \code{designed} and \code{dlsymdev}, the same code on toy inputs (110 of 110 seed-0 cells on \code{designed} checked equal), and v2's own on \code{dlfamdev}.
\item \emph{Image component: \code{p10_codes} adopted} (23:40). Both candidates qualified, and the rule takes the one with the lower mean $\log_2$ error ratio against v2 over the 8 development cells. \code{p10_codes} had $-0.078$; at \code{qd_a}, $n=60000$, 10.82\% error (limit 12.57\%), 1453 CPU s (limit 2636 s) and 0.58 of the FLOP budget. \code{p10_hier} had $-0.071$, 11.08\%, 924 s and 0.44. The builders' own numbers used Part 9's \code{p9_gates} development errors for v2 on \code{qd_dev} and Fashion-MNIST (so \code{p10_hier}'s note reports $-0.068$); the rule used v2's own runs on those cells, made for Part 10.
\item \emph{Assembly and freeze} (23:45; commit \texttt{20003f85}, 23:50). The script \repo{study/part10/assemble_v3.py} built v3 from the learner file of \code{p10_codes}, with DFG v1.3. \texttt{learner.py} has md5 \texttt{b045fafe1e24418620bf8c503aa601c6} and \texttt{dfg13.py} \texttt{69cabe5c50345cb24079198ed52b562f}. All 15 identity checks on development data gave identical predictions: v3 against \code{p10_sum} on 12 toy fits and on kr\_vs\_kp, against \code{p10_codes} on a \code{qd_dev} fit, and against v2 on optdigits. The files were made read-only.
\item \emph{Scheduling.} Neither image candidate changes any non-image code path, so v3's behaviour on toy inputs was fixed by the toy decision alone. v3's held-out toy evaluation therefore started after the toy decision, as \code{p10_sum} on \code{dlheld7}. Its results were not read before v3 was frozen, and v3 $\equiv$ \code{p10_sum} on toy inputs was checked at assembly.
\item \emph{Best-learner rule: v3} (commit \texttt{adcbfd0c}, 04:20 on 28 October; the protocol entry says 04:25). It was applied after every held-out cell of v3 and v2 was in. The mean $\log_2$ error ratio of v3 to v2 over the four held-out image cells with $n\ge10000$ is $-0.125$ (rule: below 0), and the mean $\log_2\nstar$ ratio on \code{dlheld7} is $-0.004$ (rule: at most $+0.1$).
\end{itemize}

\subsection{Disclosures and notes}\label{sec:disc-provenance-disclosures}

\paragraph{What was seen of held-out data.} The protocol and the builders' notes record three items.
\begin{itemize}
\item \emph{The \code{p10_hier} builder}, about 15 minutes after its freeze, printed the last 5 lines of the orchestrator's log of the baseline \code{mlp} on \code{qd_c} while finding out why its $n=60000$ cell was waiting for the shared lock. Those lines are the baseline MLP's \code{qd_c} test errors at $n=100$--3000. Nothing about a discrete learner on \code{qd_c} was shown, and the learner file was already frozen and read-only. \code{p10_hier} was not adopted.
\item \emph{The \code{p10_sum} builder} recorded incidental exposures: a process listing that showed the command lines of reference jobs on \code{qd_c}, and a one-line function signature \texttt{make\_dl\_heldout7\_suite} printed at the end of a source range, with no body. Nothing about \code{dlheld7}'s targets was seen.
\item \emph{The orchestrator}, before v3's freeze, read reference results on held-out data: the \code{nets_wd} training logs of f7\_00--f7\_04 on \code{dlheld7} and the MLP references' \code{qd_c} progress lines. It read no result of v2 or of a Part 10 learner on \code{dlheld7}, \code{qd_c} or \code{qd_d}. v2's \code{dlheld7} results and v3's toy results there were committed unread.
\end{itemize}

\paragraph{BLAS threading.} The image analyst found at about 02:00 on 28 October that v2's own development cells on \code{qd_dev} and Fashion-MNIST had been run without forcing single-threaded BLAS, while the candidates' cells and every reference lane force it. Thread count changes summation order, which can move nearest-neighbour ties and booster early stopping: v2 at \code{qd_dev}, $n=10000$, seed 0 gives 0.1463 multi-threaded and 0.1467 in a single-threaded analysis refit (\code{ggbt} stopped at 166 and 123 rounds). An earlier analysis refit, with unrecorded thread settings, reproduces the cell's 0.1463 exactly, so the BLAS explanation fits the evidence but was not verified. The integration rule compared both image candidates with the same v2 numbers, so the choice between them ($-0.078$ against $-0.071$) is unaffected, and both are far from the $-0.05$ threshold. Every Part 10 held-out run, of v3, v2 and the references, is single-threaded. The note was added to the protocol at 02:03 (\texttt{a2b8470e}).

\paragraph{Memory and job control.} Up to about ten jobs shared 4 cores and 16 GB. The builders ran at nice 0. Every reference and analysis process ran at nice $\ge5$, was resumable (per cell for the image evaluations, per target for the networks), and had a raised \texttt{oom\_score\_adj}, so that under memory pressure the kernel would kill it before a builder's process or v3's evaluation. Large fits ($n\ge30000$) took a shared lock. The out-of-memory kills were:
\begin{itemize}
\item the baseline \code{mlp} on \code{qd_c}, at about 20:00--20:15 during the builders' evaluations, after its $n\le10000$ cells had finished;
\item \code{p10_hier}'s \code{qd_dev} evaluation process (about 20:29, above);
\item v2's held-out image evaluation, more than once between about 22:00 and 00:30. One kill, at about 00:24, happened while two image analysis fits were running; the analyst then held its own processes at the highest kill priority, and v2's held-out cells were moved to nice 0 and out of the killer's preference, since they feed the best-learner rule;
\item reference processes while v3's $n=60000$ cells ran: v1 on \code{qd_c}, \code{mlp_xl} on \code{qd_d}, \code{mlp_pool} and \code{cnn} on \code{qd_c};
\item \code{nets_cewd0.3} and \code{nets_cewd1.0} on \code{dlheld7} at about 02:40, with 6 of 12 targets saved each (both have since completed all 12).
\end{itemize}
Every killed job was rerun with the same code and seeds, keeping its finished cells or targets. The container restart at 22:40 also stopped every job lane; the lanes were resumed the same way.

\subsection{The interpretability workflow}\label{app:interp-workflow}

\paragraph{Analysts.} Two analysts, one for the toy side and one for images, ran instrumented reruns of v3's (and, for comparison, v2's) seed-0 evaluation fits, with the learner's internals recorded (\cref{sec:interp}). The protocol fixed the analysis as analysis of the chosen learner, to be run after the held-out results were in; in practice it ran in parallel with the held-out evaluation, on v3 before the best-learner rule had chosen it. The image analysis used \code{qd_dev}, \code{qd_a}, MNIST and Fashion-MNIST and opened no result on \code{qd_c} or \code{qd_d}. The toy analysis covers every toy suite, and the six sum-structured ones in detail, \code{dlheld7} included; its \code{dlheld7} runs were made from 00:16 on 28 October, after v3's freeze. The work ran in parallel with the held-out evaluation and was committed at 03:54 on 28 October (\texttt{9214cfde}), before the best-learner rule was applied at 04:20. The rule is mechanical, and the analysis did not touch its inputs. The protocol fixed the analysis as descriptive, not as evidence that the learner is better than anything.

\paragraph{Identity of the analysed fits.} Every analysed fit is seed 0, and each $(\text{target},n)$ or $(\text{set},n)$ pair is one fit. The fits are instrumented reruns of the evaluation's own fits: the same training rows, column permutation and random streams, with the learner's internals recorded.
\begin{itemize}
\item \emph{Toy.} At all 875 cells of v3 (every toy suite) and all 477 cells of v2 (the six sum-structured suites) that have an evaluation result, the reruns' test errors equal the evaluation's seed-0 errors. Equal errors at every cell are strong evidence that the fits are the same, but not a proof. The analysis's own check (\texttt{check\_identity} in \repo{study/part10/interp/toy_interp.py}) covered 587 of the v3 cells, because v3's result files on \code{random} and \code{dlheld}--\code{dlheld3} were incomplete when it ran. \repo{paper_learner/check_toy_identity.py} repeats the same comparison on the final result files (output: \texttt{check\_toy\_identity.json}); it adds the 288 cells of those four suites, all equal.
\item \emph{Images.} The four v3 fits that have an evaluation cell (\code{qd_dev} at $n=1000$ and $10000$, \code{qd_a} and Fashion-MNIST at $n=10000$) reproduce its test error, stacking weights, \code{ggbt}'s number of rounds and every member's out-of-fold error. The fits at \code{qd_dev}, $n=3000$, and MNIST, $n=10000$, have no cell and are unchecked. No result on \code{qd_c} or \code{qd_d} was opened. v2's analysis fit at \code{qd_dev}, $n=10000$, differs from its evaluation cell by BLAS threading (\cref{sec:disc-provenance-disclosures}).
\end{itemize}

\paragraph{Verifications.} The analyses are \repo{study/part10/interp/TOY_INTERP.md} and \repo{study/part10/interp/IMAGE_INTERP.md}, with their scripts, tables and result files in the same folder; the verifications are in \repo{study/part10/interp/VERIFICATION.json}. Each side had two independent verifiers. One recomputed every number from the files with its own scripts; the other reviewed the readings against the files and the code. Where a verifier corrected a number or a reading, \cref{sec:interp} uses the corrected one. \Cref{tab:verif} counts the verdicts. ``Corrected'' means that the verifier changed a number, narrowed a statement or rejected a reading. Each verifier also listed further problems that were not tied to one claim.

\begin{table}[ht]
\centering\footnotesize\setlength{\tabcolsep}{5pt}
\begin{tabular}{llcccc}
\toprule
side & verifier & claims & confirmed & corrected & further problems listed\\
\midrule
toy & numbers & 13 & 8 & 5 & 12\\
toy & reading & 13 & 2 & 11 & 13\\
images & numbers & 11 & 10 & 1 & 12\\
images & reading & 11 & 2 & 9 & 11\\
\bottomrule
\end{tabular}
\caption{The verifications of the interpretability analyses. Source: \repo{study/part10/interp/VERIFICATION.json}.}
\label{tab:verif}
\end{table}

The corrections that change numbers in \cref{sec:interp} are these:
\begin{itemize}
\item \emph{The detection test for modulated blocks}, recomputed as DFG tests it (\cref{app:interp-mixtures}): seven gate and multiplexer targets became detectable later, so the targets detectable at $n=128$ went from 43 to 39, the median delay over the targets found from 1 doubling to 0, and Spearman's $\rho$ over all 54 from $-0.44$ to $-0.54$.
\item \emph{Counts and scopes.} The beam and forced blocks move 5 pure symmetric targets plus one gate mixture, not 6 pure targets. The pooling-set locality fractions exclude the centre pixel. The share of evidence on blank pixels is not above one half for every class; several classes are listed as exceptions.
\end{itemize}
The corrections of readings are these:
\begin{itemize}
\item the ``learned function of the block'' is the whole DFG model's conditional mean, not the block's own coefficients;
\item for fd\_10 the pair statistics are small (about 0.002) but not zero;
\item stacking-weight shares are not contributions to the prediction;
\item split-gain shares must be read against column shares;
\item the bars' straightness is a property of the construction and the data, not of the booster;
\item the regions named in the contribution maps were chosen after looking at them.
\end{itemize}
\Cref{sec:interp} uses the corrected numbers and wording; where the claims check below found a correction that had not been carried over, it was applied. \texttt{make\_interp.py} hard-codes the verifier's corrected detection points and regenerates every interpretability figure and table from the analysis files.

\subsection{This paper and reproducibility}\label{app:repro}

\paragraph{This paper.} Its sections were drafted by the orchestrating session and by separate writing agents, each from the records it names, and then compiled together. Every table and figure is generated from the result files: \texttt{make\_tables.py} and \texttt{make\_figures.py} from \repo{study/part10/LEARNER_RESULTS.json} and the result files it reads, and \texttt{make\_interp.py} from the analysis files, applying the verifiers' corrections. The plotted values are in \texttt{figures/*.json} (\texttt{fig\_toy\_nstar.json}, \texttt{fig\_images.json}, \texttt{fig\_interp.json}). \repo{study/part10/LEARNER_RESULTS.json} is written by \repo{study/part10/compare_learner.py}, whose docstring says which result file stands for v3 on each input type.

\paragraph{Which result files stand for v3.}
\begin{itemize}
\item \emph{Toy.} On toy inputs v3 runs exactly the code of the Part 10 toy candidate \code{p10_sum}; its predictions were checked identical at assembly (\repo{study/part10/hyb_portfolio3/ASSEMBLY.md}). v3's toy results are therefore the \code{p10_sum-perm} result files in \repo{study/hybrid/results/}. v2's are the \code{hyb_portfolio2-perm} files, and on \code{designed} and \code{dlsymdev} the \code{p9_sum-perm} files (the same code on toy inputs).
\item \emph{Images.} On \code{qd_c} and \code{qd_d} they are v3's own runs (\code{hyb_portfolio3}). On the development sets it is the frozen image candidate \code{p10_codes}, whose predictions on images are identical to v3's (checked identical on one \code{qd_dev} fit at assembly, $n=300$; \repo{study/part10/hyb_portfolio3/ASSEMBLY.md}).
\end{itemize}

\paragraph{Claims check.} After drafting, three independent AI checkers, each with its own instructions, checked the paper's claims on 28 October. One covered the two results sections, one the interpretability section, and one the abstract, the introduction, \cref{sec:learner,sec:protocol,sec:discussion} and this appendix. Each compared the text with the result files, the records named in it and the learner code, recomputed what it could, and listed each problem with a severity (error, misleading, minor) and a proposed correction. They found stale counts from runs that had landed after drafting, a few arithmetic and transcription errors (for example the \code{cnn}'s relative error on \code{qd_a}, and the order of HELDOUT2 and v1's freeze), several verifier corrections of the interpretability analyses that had not been carried over, statements broader than the evidence (for example ``fail because'' for a descriptive seed-0 analysis), and details of the learner description that did not match the code. Their corrections were applied before the final regeneration of every table and figure, with counts updated for the runs that landed in the meantime (\code{random} for v3, the remaining \code{dlheld7} network targets, and the image references on \code{qd_c} and \code{qd_d}). The check is not an audit of leakage, compute or reference strength of the kind made in Parts 8 and 9.

\paragraph{Editorial revision.} After the claims check, three reviewing agents (on readability, structure and layout) commented on the presentation, and an editing agent reorganised the paper: it added the comparison with familiar learners and the hypothesis class of \cref{sec:learner}, \cref{tab:headline}, the glossary and target definitions of \cref{app:names}, and moved in-sample and per-target detail to \cref{app:further,app:interp-detail}. A further agent compared every number of the revised text with the checked text, and the few new ones (the spread of the \code{dlheld7} ratios, the split of the 30 calibration wins, the \code{cnn}'s recipe) with the result files and code.

\paragraph{Regenerating the numbers.} From \texttt{gd-circuit-growth/}:
\begin{itemize}\raggedright
\item \texttt{cd study/part10 \&\& python3 compare\_part10.py} (the rules; writes \texttt{RESULTS.md});
\item \texttt{python3 compare\_learner.py} (every comparison; writes \texttt{LEARNER\_RESULTS.\{md,json\}});
\item \texttt{cd ../../paper\_learner}, then \texttt{python3 make\_tables.py}, \texttt{python3 make\_figures.py}, \texttt{python3 make\_interp.py} and \texttt{python3 check\_toy\_identity.py};
\item \texttt{latexmk -pdf main.tex}.
\end{itemize}
The result files are:
\begin{itemize}\raggedright
\item \texttt{<learner>\_\_<dataset>.json}, with \texttt{.info.json} records, in \repo{study/standard_tasks/results/};
\item \texttt{toy\_\_<learner>\_\_<suite>.json} in \repo{study/hybrid/results/}.
\end{itemize}

\paragraph{What is not in the repository.} The job drivers and queue scripts of Part 10, the builders' development scripts and logs, and the analyses' raw outputs stayed in the session scratch directory. The drivers only wrap the committed scripts, the builders' notes give the commands of their evaluations, and copies of the analysis outputs that the paper uses are in \repo{study/part10/interp/results/}.
